% ===========================================================================
% conclusion.tex — 结论
% 对应格式要求：
%   结论是全文的最终总结
%   精炼、完整、准确，不是正文中各段小结的简单重复
%   结论编排于正文之后，参考文献之前
% ===========================================================================

\chapter{结论}
\label{chap:conclusion}

本研究以水稻为研究对象，基于无人机多光谱遥感影像，结合机器学习与深度学习方法，系统开展了水稻物候期监测识别与产量估测研究。主要结论如下：

\vspace{0.3cm}

\textbf{（1）构建了基于无人机多光谱影像的水稻物候期识别方法。}
通过提取关键植被指数，结合随机森林和支持向量机算法，实现了水稻分蘖期、拔节期、抽穗期、灌浆期和成熟期的准确识别，总体分类精度达到95\%以上。其中，基于红边波段和近红外波段组合的植被指数（如NDRE）在物候期判别中表现最优。

\vspace{0.3cm}

\textbf{（2）建立了多时相无人机遥感数据融合的水稻产量估测模型。}
利用水稻全生育期的多光谱数据，融合抽穗期和灌浆期的植被指数与纹理特征构建的产量估测模型精度最高（R²~=~0.87，RMSE~=~0.45~t/ha），表明关键生育期的特征融合策略可以有效提高估测精度。

\vspace{0.3cm}

\textbf{（3）提出了基于时序植被指数动态特征的水稻早期产量估测方法。}
利用移栽至抽穗期的时序遥感数据，提取植被指数曲线动态特征参数，构建了收获前30--45天即可进行产量估测的早期预测模型（R²~=~0.78），为农业生产管理提供了及时的决策支持。

\vspace{0.3cm}

\textbf{（4）开发了基于CNN-LSTM深度学习的端到端产量估测方法。}
直接从无人机多光谱影像序列中学习时空特征，避免了人工特征提取的繁琐过程。与随机森林、支持向量机等传统机器学习方法相比，深度学习方法进一步提高了估测精度（RMSE降低约15\%），展现了深度学习在作物生长监测中的巨大潜力。

\vspace{0.5cm}

\textbf{本研究的主要创新点：}

\begin{enumerate}[itemsep=0pt,parsep=0pt]
  \item 提出了融合红边波段植被指数与机器学习的无人机遥感水稻物候期精准识别方法，显著提高了物候判别的准确性和自动化程度。
  \item 构建了基于时序植被指数动态特征的水稻早期产量估测模型，实现了收获前一个月的相对准确估产，为早期产量预测提供了新思路。
  \item 建立了基于CNN-LSTM混合模型的端到端产量估测框架，实现了从多光谱影像到产量的直接映射，为田间尺度的精准产量估测提供了智能化解决方案。
\end{enumerate}

\vspace{0.5cm}

\textbf{研究展望：}

本研究仍存在一些不足和尚待深入研究的问题。未来可在以下方面进一步开展研究：

\vspace{0.3cm}

（1）本研究仅以单一品种水稻为研究对象，未来应在多个品种、多种种植条件下进行验证，以提高模型的泛化能力。

\vspace{0.3cm}

（2）本研究主要利用多光谱数据，随着卫星遥感数据分辨率的提高和合成孔径雷达等新型传感器的发展，多源遥感数据融合将是重要的发展方向。
